% PM01 - HAI CAU CHUYEN GIAO TU KY THUAT IOP/AP (BAN THAO, KHONG PHAI DE GOC)
% Chi BODY; dung \begin{ex}...\choiceTF...\loigiai...
% Loi giai va nguon: chuyen_de/PM01_TIEN_TRINH_TU_DUY_KHONG_SO_V1.md

\begin{ex}
Một nhóm học sinh đo nhiệt dung riêng của chất lỏng bằng cùng một nhiệt lượng kế có dây đốt. Hai lần đun sử dụng cùng khối lượng chất, cùng nhiệt độ ban đầu và thời gian đun, nhưng công suất lần hai lớn hơn lần một. Nhiệt độ chất lỏng trong cả hai lần đều cao hơn nhiệt độ phòng; ở các thời điểm tương ứng sau lúc bắt đầu đun, nhiệt độ lần hai luôn cao hơn lần một. Coi nhiệt dung riêng của chất và nhiệt dung của bình không đổi, không có chuyển thể. Giả sử công suất truyền nhiệt ra môi trường tăng theo chênh lệch nhiệt độ giữa nhiệt lượng kế và phòng, với cùng quy luật cho hai lần đo.
\choiceTF
{Vì thời gian đun bằng nhau nên nhiệt lượng truyền ra môi trường trong hai lần chắc chắn bằng nhau.}
{\True Theo mô hình đã cho, nhiệt lượng truyền ra môi trường ở lần hai lớn hơn lần một.}
{\True Phương trình năng lượng phải tính cả nhiệt lượng làm nóng chất, làm nóng bình và truyền ra môi trường.}
{Nếu đã xác định chính xác nhiệt lượng hao phí mỗi lần, việc giữ nguyên khối lượng chất vẫn đủ để tìm riêng nhiệt dung riêng của chất và nhiệt dung của bình bằng hai phép đo.}
\loigiai{Gọi $Q_{{\rm hao},i}$ là nhiệt lượng truyền ra môi trường. Ta có
\[P_it=(mc+C)\Delta T_i+Q_{{\rm hao},i}.\]
Cùng thời gian không có nghĩa cùng nhiệt lượng hao phí. Vì nhiệt độ lần hai cao hơn tại mọi thời điểm tương ứng và công suất mất nhiệt tăng theo chênh lệch nhiệt độ, $Q_{{\rm hao},2}>Q_{{\rm hao},1}$. Dù biết được $Q_{{\rm hao},i}$, mỗi phép đo chỉ xác định được tổ hợp $mc+C$; không thể tách hai ẩn nếu $m$ không đổi. Đáp án: S--Đ--Đ--S.}
\end{ex}

\begin{ex}
Một phòng thí nghiệm dùng bình kín có thành cứng chứa khí lí tưởng và bộ sưởi. Cảm biến đo áp suất tuyệt đối $p$ và nhiệt độ tuyệt đối $T$; khí không rò rỉ và số phân tử khí không đổi, nhưng thể tích thực của phần chứa khí chưa được xác định. Nhóm nghiên cứu thay đổi nhiệt độ, thu được nhiều cặp số đo $p,T$ với tỉ số $p/T$ gần như không đổi trong giới hạn sai số. Một học sinh cho rằng không cần đo thể tích: chỉ cần nhiều cặp $p,T$ là xác định được số phân tử khí và chứng minh mô hình khí lí tưởng đúng trong mọi điều kiện.
\choiceTF
{\True Thành bình cứng làm thể tích khí không đổi trong quá trình đo, theo giả thiết.}
{Chỉ với các cặp số liệu $p,T$, chưa có thông tin về thể tích, vẫn xác định duy nhất số phân tử khí.}
{Kết quả $p/T$ gần như không đổi trong khoảng khảo sát đủ để chứng minh khí tuân theo định luật khí lí tưởng ở mọi điều kiện.}
{\True Nếu xác định thêm thể tích thực của khí, có thể dùng $p,T,V$ và mô hình khí lí tưởng để tính số phân tử khí.}
\loigiai{Với mô hình khí lí tưởng, $pV=Nk_{\rm B}T$, nên $\frac pT=\frac{Nk_{\rm B}}V$. Khi $V$ chưa biết, đo nhiều cặp $p,T$ trong cùng bình chỉ xác định tổ hợp $N/V$. Thành bình cứng chỉ cho biết $V$ không đổi, không cung cấp giá trị $V$. Kết quả phù hợp giả thuyết trong khoảng khảo sát không chứng minh tính đúng đắn trong mọi điều kiện. Đáp án: Đ--S--S--Đ.}
\end{ex}
